\subsection{Audio Generation}
\label{sec:audio_gen}

The primary objective of audio generation in this work is to produce synchronized soundtracks for video clips, formulated as a video-to-audio (V2A) generation framework. 
The generated soundtracks consist of ambient sound and background music, explicitly excluding speech or vocal elements. 
In contrast to text-to-audio models ~\citep{audioldm, audioldm2}, the ambient sound generated by video-to-audio models must be temporally aligned with the visual content of the video, while the accompanying music should accurately reflect the emotional tone and contextual setting of the video, as would naturally be expected. Furthermore, to enhance user control over sound design, our framework integrates both video and textual prompts, enabling users to specify on-screen or off-screen sounds and define the style and presence of background music.

\subsubsection{Model Design}
Consistent with our video generation framework, our video-to-audio model also utilizes a diffusion transformer (DiT) ~\citep{dit} based on the flow-matching~\citep{flow-matching} to model the denoising diffusion process in the audio domain. The detailed pipeline of our V2A model is illustrated in Fig.~\ref{fig:v2av2-pipeline}. 

\begin{figure}[t]  
    \centering  
    \includegraphics[width=0.99\textwidth]{./figures/audio/v2av2-pipeline.pdf}  
    \caption{Framework of video-to-audio generation. Our model processes both a video chunk and its associated textual description as dual inputs to synthesize high-quality, semantically coherent audio.} 
    \label{fig:v2av2-pipeline} % 
\end{figure}

\begin{figure}[htbp]  
    \centering  
    \includegraphics[width=0.99\textwidth]{./figures/audio/v2av2-goodcases-pickup.pdf}  
    \caption{Audio generation samples of MMAudio and ours.} 
    \label{fig:v2av2-goodcases-pickup} % 
\end{figure}

\textbf{Audio autoencoder}. 
A widely adopted audio compression approach involves transforming raw waveforms into mel-spectrograms, followed by encoding them into $H^{a}\times W^{a}\times C^{a}$ latent features using an image-based variational autoencoder, as demonstrated in ~\citep{audioldm, audioldm2}. However, in the context of a diffusion transformer backbone, these latent features must be partitioned into patches and reshaped into a tensor of size $(H^a\cdot W^a)\times C^a$. This patching-and-reshaping process can disrupt the temporal alignment with the corresponding video content, posing a significant challenge for synchronization. Considering this limitation, we train a compression model that operates directly on raw waveforms. Our 1D-VAE generates latent features of size $T^{a}\times C^{a}$, where $T^a$ represents the sequence length along the time axis, preserving explicit temporal information essential for accurate synchronization.

\textbf{Video and text encoder}. To achieve frame-accurate synchronization between synthesized audio and visual sequences, we establish temporal coherence through multimodal feature fusion. Our architecture first extracts frame-wise visual embeddings using a CLIP~\cite{radford2021learning} model, then performs temporal rate adaptation through feature replication to match the audio feature's sampling rate, following~\citep{polyak2024moviegencastmedia}. The synchronized visual and acoustic features undergo dimension transformation via linear projection layers to align with DiT's latent space, followed by element-wise summation for multimodal fusion. For linguistic understanding across languages, we leverage the pre-trained umT5 model~\citep{chung2023unimax} as our text encoder, capitalizing on its cross-lingual representation capabilities through frozen parameters.

\textbf{Data.} The training data for our V2A model is derived from the video generation dataset through a rigorous filtering process. We systematically remove videos lacking soundtracks or containing speech/vocal music, resulting in a refined subset of $\mathcal{O}(1)$ thousand hours. To enhance multimodal representation, we employ a comprehensive captioning strategy: dense video descriptions are complemented with audio-specific captions generated by Qwen2-audio~\citep{qwen2-audio}. These audio captions are categorized into ambient sounds and musical compositions, with the latter characterized by style, rhythm, melody, and instrumentation. The final structured caption integrates three components: (1) dense video description, (2) ambient sound characterization, and (3) background music analysis, providing a unified multimodal representation for training.

\textbf{Implementation details.} Our V2A architecture generates high-fidelity stereo audio with a maximum duration of 12 seconds at a 44.1 kHz sampling rate. The model processes multimodal inputs through dedicated encoders: umT5-XXL produces 4096-dimensional text embeddings, while the CLIP extracts 1024-dimensional visual embeddings per frame. The audio and video features are projected into a unified 1536-dimensional latent space within the DiT backbone. For temporal alignment, we downsample input videos and yield precisely 48 frames for 12-second clips, ensuring frame-accurate synchronization between visual and audio modalities. The system processes latent sequences of length 256 in one batch. To enhance the model's capability of generating audio solely from visual cues, we implement a random masking strategy during training where ambient sound and music captions are selectively omitted with a predefined probability. This conditioning mechanism forces the model to establish robust cross-modal associations between visual content and corresponding audio patterns, while maintaining the flexibility to utilize textual descriptions when available.


\subsubsection{Evaluation}
We provide qualitative evaluations of audio generation in Fig.~\ref{fig:v2av2-goodcases-pickup}, where we compare our method with the recently released open-sourced approach MMAudio~\citep{mmaudio}. This comparative analysis is particularly meaningful as MMAudio represents the most similar and competitive baseline to our method, demonstrating impressive performance in audio generation tasks. The videos are generated by our text-to-video model with a duration of 5 seconds.

Our V2A model demonstrates superior performance in several key aspects of audio generation, as evidenced by the comparative results. Specifically, the model exhibits enhanced long-term consistency, particularly noticeable in the `pouring coffee' scenario (the first case in Fig.~\ref{fig:v2av2-goodcases-pickup}). In the second case of `typing', our approach generates significantly cleaner audio output compared to the noisier results produced by MMAudio. Furthermore, our method excels in synthesizing rhythmic sound patterns, as demonstrated in the `walking', `boxing', and `clip-clop' examples (the last three cases in Fig.~\ref{fig:v2av2-goodcases-pickup}), where it achieves more natural and coherent audio generation.

\textbf{Limitations.} Our method exhibits limitations in generating human vocal sounds, including but not limited to laughter, crying, and speech. This limitation is primarily attributed to our data preparation process, where speech-related vocal data is deliberately excluded from the training dataset. The MMAudio model's ability to produce random speech-like sounds stems from its retention of speech-related data in its training corpus. In our future work, we plan to incorporate speech generation capabilities to address this limitation.

